\documentclass[12pt]{article}

\usepackage{amsmath,amssymb,amsthm}
\usepackage{hyperref}
\usepackage{microtype}
\usepackage[margin=1.1in]{geometry}

% ----------------------------------------------------------------
% Theorem environments
% ----------------------------------------------------------------
\newtheorem{theorem}{Theorem}
\newtheorem{lemma}[theorem]{Lemma}
\newtheorem{corollary}[theorem]{Corollary}
\newtheorem{definition}[theorem]{Definition}
\newtheorem{remark}[theorem]{Remark}

% ----------------------------------------------------------------
% Notation
% ----------------------------------------------------------------
\newcommand{\R}{\mathbb{R}}
\newcommand{\Z}{\mathbb{Z}}
\newcommand{\E}{\mathrm{E}}         % bisector energy
\newcommand{\Erot}{\mathrm{E}_{\mathrm{rot}}} % rotation energy
\newcommand{\bis}[2]{\beta(#1,#2)}  % perpendicular bisector
\newcommand{\ip}[2]{\langle #1,#2\rangle}
\newcommand{\norm}[1]{\|#1\|}
\DeclareMathOperator{\card}{card}

% ----------------------------------------------------------------
\begin{document}

\title{The minimum bisector energy of a planar point set}
\author{Adam McKenna}
\date{}
\maketitle

% ----------------------------------------------------------------
\begin{abstract}
The \emph{bisector energy} $\E(P)$ of a finite point set $P\subset\R^2$, introduced
by Lund, Sheffer, and de~Zeeuw, counts ordered quadruples of distinct point pairs
sharing a perpendicular bisector.
We prove that $\E(P)\ge 2n(n-1)$ for every $n$-point planar set,
with equality if and only if every two distinct point pairs have distinct
perpendicular bisectors (\emph{bisector-injective}).
The bound is tight: for every finite set $G$ in any Euclidean space with no
three collinear points, a single linear projection $T\colon\R^d\to\R^2$ exists
such that the image $T(G)$ is bisector-injective, hence attains the floor
$\E(T(G))=2n(n-1)$.
The projected set also has zero rotational energy and its distance multiset
is in bijection with the $\pm$-difference classes of $G$.
The proof is verified in Lean~4 against standard Mathlib axioms only
(propext, Classical.choice, Quot.sound).
\end{abstract}

% ----------------------------------------------------------------
\section{Introduction}

For a finite set $P\subset\R^2$ with $|P|=n$, define the \emph{bisector energy}
\[
  \E(P) \;=\;
  \bigl|\{(a,b,c,e)\in P^4 : a\ne b,\; c\ne e,\;
         \bis{a}{b} = \bis{c}{e}\}\bigr|,
\]
where $\bis{p}{q} = \{x\in\R^2 : |x-p|=|x-q|\}$ is the perpendicular bisector
of the pair $(p,q)$.
This quantity was introduced by Lund, Sheffer, and de~Zeeuw~\cite{LSdZ16},
who proved the upper bound
\[
  \E(P) \;=\; O\!\left(M(n)^{2/5}\,n^{12/5+\varepsilon}+M(n)\,n^2\right),
\]
where $M(n)$ is the maximum number of points of $P$ on any line or circle,
and applied it to obtain structural results about point sets with few distinct
distances.

The present note addresses the complementary question: \emph{how small can
$\E(P)$ be?}

Every unordered pair $\{a,b\}\subseteq P$ ($a\ne b$) contributes exactly four
ordered quadruples to $\E(P)$ by taking both orderings of each pair:
$(a,b,a,b)$, $(a,b,b,a)$, $(b,a,a,b)$, $(b,a,b,a)$.
These \emph{trivial quadruples} account for $4\binom{n}{2}=2n(n-1)$ contributions,
so $\E(P)\ge 2n(n-1)$ for every $P$.
A nontrivial contribution arises precisely when two distinct unordered pairs
share a perpendicular bisector.
We say $P$ is \emph{bisector-injective} if no such coincidence occurs,
i.e.\ the map $\{a,b\}\mapsto \bis{a}{b}$ is injective on unordered pairs.

The main theorem shows that the floor $2n(n-1)$ is attained.

\begin{theorem}[Near Enemy Theorem for Bisector Energy]
\label{thm:main}
Let $G$ be a finite set in a Euclidean space $\R^d$ with no three collinear
points, and let $n=|G|$.  There exists a linear map
$T\colon\R^d\to\R^2$ such that:
\begin{enumerate}
  \item[\rm(a)] $T$ is injective on $G$;
  \item[\rm(b)] $\E(T(G)) = 2n(n-1)$ \emph{(bisector-energy floor)};
  \item[\rm(c)] $\E(T(G))\le \E(P')$ for every $n$-point planar set $P'$
        \emph{(absolute minimality)};
  \item[\rm(d)] $T(G)$ has no three collinear points and no four concyclic
        points \emph{(full general position)};
  \item[\rm(e)] $\Erot(T(G))=0$ \emph{(zero rotational energy)};
  \item[\rm(f)] the distance multiset of $T(G)$ is in bijection with
        the $\pm$-difference classes of $G$
        \emph{(distance transport)}.
\end{enumerate}
\end{theorem}

The EFPR construction---the lattice-sphere slice $G=\{x\in[-h,h]^d\cap\Z^d :
\norm{x}^2=R\}$ of Erd\H{o}s, F\"{u}redi, Pach, and Ruzsa~\cite{EFPR93}---satisfies
the no-three-collinear hypothesis (a line meets a sphere in at most two points),
so the theorem applies.  The projected EFPR family is the canonical minimizer.

\medskip
\noindent\textbf{Notation.}
We write $\ip{\cdot}{\cdot}$ for the standard inner product on $\R^d$,
$\norm{\cdot}$ for the induced norm, and $|S|$ for the cardinality of a
finite set $S$.
For a linear map $T\colon\R^d\to\R^2$ we write $T_1,T_2\in\R^d$ for its
two row vectors, so $T(v)=({\ip{T_1}{v}},{\ip{T_2}{v}})$.
The \emph{$\pm$-difference classes} of $G$ are the equivalence classes of
$G{-}G\setminus\{0\}:=\{a-b:a,b\in G,\,a\ne b\}$ under $u\sim v$ iff $u=\pm v$.

% ----------------------------------------------------------------
\section{Definitions and the lower bound}

\begin{definition}
A finite set $P\subset\R^2$ is \emph{bisector-injective} if
$\bis{a}{b}\ne\bis{c}{e}$ for every two distinct unordered pairs
$\{a,b\},\{c,e\}\subseteq P$ with $a\ne b$ and $c\ne e$.
\end{definition}

\begin{lemma}[Floor and equality]
\label{lem:floor}
For every $n$-point set $P\subset\R^2$:
\[
  \E(P) \;\ge\; 2n(n-1),
\]
with equality if and only if $P$ is bisector-injective.
\end{lemma}

\begin{proof}
For each ordered pair $(a,b)$ with $a\ne b$, the quadruples $(a,b,a,b)$ and
$(a,b,b,a)$ both lie in the defining set of $\E(P)$ since $\bis{a}{b}=\bis{b}{a}$.
As $(a,b)$ ranges over the $n(n-1)$ ordered pairs with $a\ne b$, these $2n(n-1)$
quadruples are pairwise distinct (the first ordered pair $(a,b)$ and the swap bit
together uniquely determine the quadruple), so $\E(P)\ge 2n(n-1)$.

If $P$ is bisector-injective then every quadruple counted by $\E(P)$ has
$\bis{a}{b}=\bis{c}{e}$, hence $\{a,b\}=\{c,e\}$ as unordered pairs,
so the quadruple is trivial.  Thus $\E(P)=2n(n-1)$.

Conversely, if $P$ is not bisector-injective there exist distinct unordered
pairs $\{a,b\}\ne\{c,e\}$ with $\bis{a}{b}=\bis{c}{e}$.
The four ordered quadruples $(a,b,c,e)$, $(a,b,e,c)$, $(b,a,c,e)$, $(b,a,e,c)$
are all counted by $\E(P)$ and are nontrivial, giving $\E(P)>2n(n-1)$.
\end{proof}

% ----------------------------------------------------------------
\section{The key criterion: shared bisectors upstairs}

The proof of Theorem~\ref{thm:main} uses a criterion that translates a shared
bisector downstairs into algebraic conditions on the preimages upstairs.

\begin{lemma}[Shared-bisector conditions]
\label{lem:shared-bis}
Let $p,q,p',q'\in\R^2$.  If $\bis{p}{q}=\bis{p'}{q'}$ then:
\begin{enumerate}
  \item[\rm(i)] $q'-p'=t(q-p)$ for some $t\in\R$ \emph{(parallel differences)};
  \item[\rm(ii)] $\ip{p+q-(p'+q')}{q-p}=0$
        \emph{(midpoint difference orthogonal to direction)}.
\end{enumerate}
\end{lemma}

\begin{proof}
Two perpendicular bisectors are equal as affine subspaces if and only if they
have the same direction and share a point.
The direction of $\bis{p}{q}$ is the orthogonal complement of $\R(q-p)$;
equal directions give (i).
Both midpoints $\frac{p+q}{2}$ and $\frac{p'+q'}{2}$ lie on the common bisector,
so their difference lies in its direction subspace, giving (ii).
\end{proof}

\begin{lemma}[Upstairs line rigidity]
\label{lem:rigid}
Let $G\subset\R^d$ have no three collinear points, and let $a,b,c,e\in G$ with
$a\ne b$ and $c\ne e$.  If
\begin{enumerate}
  \item[\rm(i)] $e-c = t(b-a)$ for some $t\in\R$, and
  \item[\rm(ii)] $\ip{a+b-(c+e)}{b-a}=0$,
\end{enumerate}
then $\{a,b\}=\{c,e\}$.
\end{lemma}

\begin{proof}
From~(i), all four points lie on the affine line $\{a+s(b-a):s\in\R\}$.
Write $c=a+s(b-a)$ and $e=a+(s+t)(b-a)$ for some $s\in\R$.
Substituting into~(ii):
\[
  \ip{a+b-(c+e)}{b-a}
  \;=\; \ip{(1-2s-t)(b-a)}{b-a}
  \;=\; (1-2s-t)\norm{b-a}^2 \;=\; 0.
\]
Since $a\ne b$, we get $t=1-2s$, hence $e=a+(1-s)(b-a)$.
If $s\notin\{0,1\}$ then $a$, $b$, $c=a+s(b-a)$ are three distinct collinear
points, contradicting the no-three-collinear hypothesis.
Hence $s\in\{0,1\}$: if $s=0$ then $c=a$ and $e=b$; if $s=1$ then $c=b$ and $e=a$.
In either case $\{c,e\}=\{a,b\}$.
\end{proof}

\begin{remark}
The sphere-slice case (where $G$ lies on a sphere) gives a cleaner proof of
Lemma~\ref{lem:rigid}: chords of a sphere with the same midpoint have equal
length (the parallelogram law expresses $\norm{a-b}^2$ in terms of $\norm{a+b}/2$
and the radius, both fixed by the common midpoint and sphere).
Two parallel chords with equal length and the same midpoint then coincide:
the midpoint pins the foot of the perpendicular from the center, and equal
length pins the two endpoints symmetrically about it.
The no-three-collinear hypothesis makes the sphere-slice hypothesis unnecessary
for the headline theorem, since a sphere already has no three collinear points.
\end{remark}

% ----------------------------------------------------------------
\section{Generic projections}

\begin{definition}
A linear map $T\colon\R^d\to\R^2$ is \emph{generic for $G$} if:
\begin{enumerate}
  \item[(inj)] $T(a)\ne T(b)$ for all distinct $a,b\in G$;
  \item[(avoid)] for every $a,b,c,e\in G$ with $a\ne b$, $c\ne e$, and
        $\{a,b\}\ne\{c,e\}$: it is not the case that both
        $T(e)-T(c)=t\cdot(T(b)-T(a))$ for some $t\in\R$ \emph{and}
        $\ip{T(a)+T(b)-(T(c)+T(e))}{T(b)-T(a)}=0$;
  \item[(sep)] for every $a,b,c,e\in G$ with $a-b\ne\pm(c-e)$:
        $\norm{T(a-b)}^2\ne\norm{T(c-e)}^2$.
\end{enumerate}
\end{definition}

\begin{lemma}[Generic projection implies bisector-injectivity]
\label{lem:gen-implies-inj}
If $T$ is generic for $G$ then $T(G)$ is bisector-injective.
\end{lemma}

\begin{proof}
Suppose $T(a),T(b),T(c),T(e)\in T(G)$ satisfy
$\bis{T(a)}{T(b)}=\bis{T(c)}{T(e)}$ with $T(a)\ne T(b)$ and $T(c)\ne T(e)$.
By injectivity of $T$ on $G$, we have $a\ne b$ and $c\ne e$.
By Lemma~\ref{lem:shared-bis}, conditions (i) and (ii) hold downstairs.
Linearity of $T$ lifts them: $T(e-c)=tT(b-a)$ and $\ip{T(a+b-(c+e))}{T(b-a)}=0$.
The avoid condition then forces $\{a,b\}=\{c,e\}$, hence $\{T(a),T(b)\}=\{T(c),T(e)\}$.
\end{proof}

\begin{lemma}[Existence of a generic projection]
\label{lem:exists-gen}
If $G\subset\R^d$ is a finite set with no three collinear points, a generic
projection $T\colon\R^d\to\R^2$ for $G$ exists.
\end{lemma}

\begin{proof}
We view the rows $T_1,T_2\in\R^d$ as $2d$ free real variables.
For each off-pair quadruple $(a,b,c,e)$ (with $a\ne b$, $c\ne e$, $\{a,b\}\ne\{c,e\}$),
define two polynomials in $(T_1,T_2)$:
\begin{align*}
  D(a,b,c,e) &= \ip{T_1}{b-a}\ip{T_2}{e-c} - \ip{T_2}{b-a}\ip{T_1}{e-c}, \\
  H(a,b,c,e) &= \ip{T_1}{a+b-(c+e)}\ip{T_1}{b-a}
                +\ip{T_2}{a+b-(c+e)}\ip{T_2}{b-a}.
\end{align*}
If both $D=0$ and $H=0$ at some $T$, then condition~(avoid) fails,
and the proof of Lemma~\ref{lem:gen-implies-inj} would give $\{a,b\}=\{c,e\}$,
contradicting the off-pair assumption.
Hence for each off-pair quadruple, $D^2+H^2$ is a nonzero polynomial in $2d$ variables;
its real zero-set has measure zero.
Similarly, for each pair of distinct $\pm$-difference classes $[u]\ne[v]$ in $G{-}G\setminus\{0\}$,
the polynomial $\sum_k\ip{T_k}{u-v}\ip{T_k}{u+v}$ (encoding condition~(sep)) is nonzero:
since $u\ne\pm v$, both $u-v\ne 0$ and $u+v\ne 0$, so neither linear factor vanishes
identically, and their inner product $\norm{Tu}^2-\norm{Tv}^2$ is a nonzero degree-2
polynomial. Likewise the injectivity condition~(inj) cuts out a measure-zero set.

The set of $(T_1,T_2)\in\R^{2d}$ satisfying all three conditions is the complement
of finitely many measure-zero algebraic sets, hence nonempty (in fact dense)~\cite{ST12}.
\end{proof}

% ----------------------------------------------------------------
\section{Proof of Theorem~\ref{thm:main}}

\begin{proof}[Proof of Theorem~\ref{thm:main}]
Parts (a)--(c) combine Lemmas~\ref{lem:floor}, \ref{lem:gen-implies-inj},
and~\ref{lem:exists-gen}: take $T$ generic for $G$, which exists by
Lemma~\ref{lem:exists-gen}.  Then $T$ is injective on $G$ (a), $T(G)$ is
bisector-injective so $\E(T(G))=2n(n-1)$ by Lemma~\ref{lem:floor} (b),
and $2n(n-1)$ is the global minimum by the lower bound of
Lemma~\ref{lem:floor} (c).

For part (d), one adds two more avoidance conditions to the definition of generic:
no collinear triple and no concyclic quadruple in $T(G)$.
For collinear triples: $T(a),T(b),T(c)$ are collinear iff
$D_{abc}(T)=\ip{T_1}{b-a}\ip{T_2}{c-a}-\ip{T_2}{b-a}\ip{T_1}{c-a}=0$.
Since $a,b,c$ are not collinear in $G$, the vectors $b-a$ and $c-a$ are linearly
independent, so $D_{abc}$ is a nonzero polynomial.
For concyclic quadruples: the $4\times 4$ circle determinant is likewise a nonzero
polynomial in $(T_1,T_2)$, provided no three of the four upstairs points are
collinear---a condition guaranteed by the no-three-collinear hypothesis on $G$.
The same measure-zero argument applies.

For part (e), write $\Erot(T(G))$ as the count of ordered quadruples
$(p,q,r,s)\in T(G)^4$ with $p\ne q$, $r\ne s$, $|p-q|=|r-s|$, and
$p-q\notin\{r-s,-(r-s)\}$.  The condition $|p-q|=|r-s|$ is equivalent to
$\norm{T(a-b)}=\norm{T(c-e)}$, which expands as
$\norm{T(a-b)}^2 - \norm{T(c-e)}^2
 = \sum_k \ip{T_k}{(a-b)-(c-e)}\ip{T_k}{(a-b)+(c-e)} = 0$.
Condition~(sep) forces $a-b=\pm(c-e)$,
i.e.\ $p-q=\pm(r-s)$, so the proper-rotation channel is empty and
$\Erot(T(G))=0$.

For part (f), the same separation argument shows that the map
$T(G)-T(G)\setminus\{0\}\to G-G\setminus\{0\}$ preserves $\pm$-classes,
giving a bijection between the sets of distinct distances.
\end{proof}

% ----------------------------------------------------------------
\section{Corollaries}

\begin{corollary}[Sphere-slice form]
\label{cor:sphere}
Let $G$ be a finite subset of a sphere in $\R^d$, $|G|=n$.  Then a linear map
$T\colon\R^d\to\R^2$ exists satisfying {\rm(a)}--{\rm(f)} of
Theorem~\ref{thm:main}.
\end{corollary}

\begin{proof}
A line meets a sphere in at most two points, so $G$ has no three collinear
points.  Apply Theorem~\ref{thm:main}.
\end{proof}

\begin{corollary}[EFPR minimizer]
The EFPR lattice-sphere-slice family $G = \{x\in[-h,h]^d\cap\Z^d :
\norm{x}^2=R\}$~\cite{EFPR93} admits a planar projection that attains the
bisector-energy floor.  In particular, the minimum bisector energy of an
$n$-point planar set is exactly $2n(n-1)$.
\end{corollary}

\begin{proof}
Apply Corollary~\ref{cor:sphere} to the EFPR lattice points on a sphere.
\end{proof}

\begin{remark}[Relation to LSdZ]
Lund, Sheffer, and de~Zeeuw~\cite{LSdZ16} remark in passing that
$\E(P) = \Omega(n^2)$, but do not extract the constant $2n(n-1)$, state the
floor, or characterize equality.
The minimization direction (Lemma~\ref{lem:floor}), the equality condition
(bisector-injectivity), and the identification of the EFPR projection as
a minimizer (Corollary~\ref{cor:sphere}) appear to be new.
\end{remark}

\begin{remark}[Rotation energy]
Part~(e) says the near-enemy projection has zero rotational energy in the
sense of the Elekes--Sharir / Guth--Katz framework~\cite{ES11,GK15}: every
congruent quadruple $(a,b,c,e)$ in $T(G)$ is realized by a translation
($a-b=c-e$) or half-turn ($a-b=-(c-e)$), never by a proper rotation.
This is the strongest possible statement: a set with any two congruent pairs
at a proper-rotation angle apart contributes to the rotation channel.
\end{remark}

% ----------------------------------------------------------------
\section{Verification}

All results in this paper are verified in Lean~4 using Mathlib.
The Lean source is
\texttt{NearEnemyTheorem.lean} in the \texttt{NearEnemy} namespace
of the \texttt{lean-formalizations} repository.
The headline Lean theorem corresponding to Theorem~\ref{thm:main} is
\begin{center}
\texttt{nearEnemy\_noThreeCollinear\_exists\_bisectorEnergy\_minimal}\\
\texttt{\_image\_generalPosition\_distanceTransport}.
\end{center}
Running \texttt{\#print axioms} on this theorem confirms that the proof
depends only on \texttt{propext}, \texttt{Classical.choice}, and
\texttt{Quot.sound} --- the standard Lean/Mathlib axioms.
No \texttt{sorry} appears anywhere in the proof chain.

% ----------------------------------------------------------------
\begin{thebibliography}{9}

\bibitem{EFPR93}
P.~Erd\H{o}s, Z.~F\"{u}redi, J.~Pach, and I.~Z.~Ruzsa,
\emph{The grid revisited},
Discrete Math.\ \textbf{111} (1993), no.\ 1--3, 189--196.

\bibitem{ES11}
G.~Elekes and M.~Sharir,
\emph{Incidences in three dimensions and distinct distances in the plane},
Combin.\ Probab.\ Comput.\ \textbf{20} (2011), no.\ 4, 571--608.

\bibitem{GK15}
L.~Guth and N.~H.~Katz,
\emph{On the Erd\H{o}s distinct distances problem in the plane},
Ann.\ of Math.\ (2)\ \textbf{181} (2015), no.\ 1, 155--190.

\bibitem{LSdZ16}
B.~Lund, A.~Sheffer, and F.~de~Zeeuw,
\emph{Bisector energy and few distinct distances},
Discrete Comput.\ Geom.\ \textbf{56} (2016), no.\ 2, 337--356.
arXiv:1411.6868.

\bibitem{ST12}
J.~Solymosi and T.~Tao,
\emph{An incidence theorem in higher dimensions},
Discrete Comput.\ Geom.\ \textbf{48} (2012), no.\ 2, 255--280.
arXiv:1103.2926.

\end{thebibliography}

\end{document}
